\section{Related Works}

In recent years, hardware intermediate representations design has become a popular topic. An appropriate IR can give sufficient expressivity and avoid redundant engineering efforts. Several existing hardware IRs are introduced next.

\paragraph{Domain-specific languages} DSLs provide a higher level of abstraction that is natural to the domain, assisting in the generation of hardware with expert knowledge. Raghu Prabhakar et al.~\cite{parallel_pattern} adopt parallel patterns like \emph{map} and \emph{fold} to express high-level computation. Aetherling~\cite{aetherling} aims at generating streaming accelerators by applying transformations to the proposed data-parallel IR. Nithin George et al.~\cite{system_dsl} generate hardware systems from applications written in a machine-learning DSL. HeteroCL~\cite{heterocl} provides an abstraction that decouples algorithm description and architecture specification, and develops a compilation flow to heterogeneous computing platforms. Apart from domain-specific optimizations, extra compilation to the RTL program is required, which makes it difficult for researchers to generate hardware quickly.

% \hl{can we merge the IRs into one paragraph?}

% \paragraph{IR for HDL} FIRRTL~\cite{firrtl} that obeys AST format is the intermediate representation in Chisel. It is a multi-level IR in which high-level abstraction is lowered by simplifying and reducing some high-level semantics like some object-oriented features. LLHD\cite{LLHD_Schuiki} is a three-level IR that includes behavioral, structural, and netlist levels of hardware synthesis. Different level in LLHD aims at different applications including simulation, verification, and logic synthesis. Hector uses higher-level abstraction to combine different synthesis methods and avoid redundant implementations, whereas these IRs are still at a low-level abstraction with no high-level information.

% \paragraph{IR for HLS} Many HLS tools~\cite{VivadoHls, legup, Dynamatic_Josipovi} still adopt LLVM IR as their internal representation, which is a pure software IR that doesn't contain any hardware information. $\mu$IR~\cite{uIR_Sharifian} uses dynamic scheduling based on task-level parallelism and briefly presents a task-level representation. Wu et al.~\cite{hierarchy_cdfg} propose a hierarchical CDFG representation for hardware and software co-design. SynASM~\cite{synASM} proposes a timed CDFG that is extended with some timing information such as sequential and parallel, and transforms it directly to VHDL using its own HLS process. AHIR~\cite{ahir} proposes a low-level abstraction that decouples the datapath and control path. All of these IRs only focus on one single abstraction level which loses the opportunity of interfering with HLS at different levels. Furthermore, these representations aren't able to represent hardware with both static and dynamic behavior.

% \paragraph{Hardware infrastructure} Calyx\cite{Calyx_Rachit} provides software-like control flow primitives such as seq, par, and while, that are then lowered to the finite state machine (FSM). It's not suitable for high-level synthesis or hardware with complex pipelines since it lacks pipeline and task-level semantics, resulting in poor expressivity. Moreover, Calyx only supports FSM-controlled hardware, so hardware with dynamic behavior like streaming accelerator is not supported. The goal of CIRCT~\cite{circt} is to construct a reusable and modular infrastructure for the entire hardware generation including high-level synthesis and logic synthesis. This project is still in progress and absorbs existing IR designs like Calyx and LLHD. In contrast, Hector concentrates on a higher abstraction than RTL, which can be regarded as the front-end of hardware generation.


\paragraph{IR for hardware} FIRRTL~\cite{firrtl} that obeys AST format is the intermediate representation in Chisel. LLHD~\cite{LLHD_Schuiki} is a three-level IR that aims at different applications including simulation, verification, and logic synthesis. $\mu$IR~\cite{uIR_Sharifian} uses dynamic scheduling based on task-level parallelism and briefly presents a task-level representation. Wu et al. and synASM~\cite{hierarchy_cdfg, synASM} propose a CDFG represention for hardware and software. AHIR~\cite{ahir} proposes a low-level abstraction that decouples the datapath and control path. Calyx\cite{Calyx_Rachit} provides software-like control flow primitives such as seq, par, and while to describe hardware. The goal of CIRCT~\cite{circt} is to construct a reusable and modular infrastructure for the entire hardware generation including high-level synthesis and logic synthesis. This project is still in progress and absorbs existing IR designs like Calyx and LLHD.